\documentclass[11pt,a4paper]{ctexart}
\usepackage[a4paper,margin=24mm]{geometry}
\usepackage{amsmath,amssymb,booktabs,longtable,tabularx,array}
\usepackage{graphicx,pdfpages,xcolor,hyperref}
\usepackage{enumitem}
\usepackage{caption}

\setCJKmainfont{Noto Serif CJK SC}
\setCJKsansfont{Noto Sans CJK SC}
\setmainfont{TeX Gyre Termes}
\setsansfont{TeX Gyre Heros}
\hypersetup{colorlinks=true,linkcolor=blue!55!black,urlcolor=blue!55!black}
\setlength{\parindent}{2em}
\setlength{\parskip}{0.35em}
\setlength{\emergencystretch}{3em}
\renewcommand{\arraystretch}{1.2}
\captionsetup{font=small,labelfont=bf}

\newcommand{\transnote}[1]{%
  \par\noindent\colorbox{yellow!16}{%
    \parbox{\dimexpr\linewidth-2\fboxsep}{\textbf{译者核对提示：}#1}}\par}

\title{EPHIC 模型：通用 SPICE 光子模型\\
用于闭环电子--光子协同仿真\\
\large 中文学习译本（非官方）}
\author{原作者：Da Ming、Yuhang Wang、Zhicheng Wang、Ken Xingze Wang\\
Ciyuan Qiu、Min Tan\\
翻译整理：项目学习稿}
\date{原文发表于 \textit{IEEE Transactions on Circuits and Systems I: Regular Papers}, 2024, 71(4), 1819--1831\\
DOI: \href{https://doi.org/10.1109/TCSI.2024.3353459}{10.1109/TCSI.2024.3353459}}

\begin{document}
\maketitle

\begin{center}
\small
本文件仅用于组内学习与术语理解，不是 IEEE 或原作者发布的官方中文版本。\par
公式、图号、表号和参考文献编号尽量与原文保持一致；原文页面完整附于文末，便于逐项核对。
\end{center}

\tableofcontents
\clearpage

\begin{abstract}
本文提出了用于闭环电子--光子协同仿真的通用 SPICE 光子模型。SPICE 模型原本面向电子器件，而电子器件与光子器件的物理机制显著不同。已有 SPICE 光子模型大多面向特定器件，不能表达多维信号等重要特征。因此，如何用 SPICE 模型表示具有关键特征的通用光子器件，仍是一个核心难题。

为解决上述问题，本文将信号的物理含义与其数值结果分离，并利用 SPICE 原生基元，以数值等效方式构建光子模型。波长、偏振等数值等效电子信号的物理含义由设计者解释。本文模型既包含现有 SPICE 模型的波长相关性和双向传输等功能，也加入了偏振、光学 Kerr 效应、双光子吸收、自由载流子吸收和自由载流子色散等新功能，并首次实现闭环偏振仿真。模型可用 DC 扫描快速完成频域仿真，不必采用逐频点瞬态仿真或基于频率啁啾的方法。

所提出的 SPICE 光子模型与 Verilog-A 模型和 Lumerical 模型的仿真结果高度一致。由于编译时间更短、兼容性更好，它与 SPICE 电子模型协同仿真时比 Verilog-A 光子模型更高效。这些通用模型使电子--光子融合电路的设计接近传统集成电路设计的便利程度，为电子与光子在电路层面的融合奠定基础。
\end{abstract}

\transnote{摘要中的“首次实现”、模型一致性和效率提升均为原文作者报告的结论，本译本未对其首创性或全部性能数据进行独立复现。}

\noindent\textbf{关键词：}建模；电子--光子融合；闭环；SPICE；电子--光子协同仿真；微环；马赫--曾德尔干涉仪；偏振。

\section{引言}

电子--光子融合（electronic-photonic convergence, EPC）是后摩尔时代的重要研究方向，目前正从器件集成走向电路设计~[1,2]。EPC 可以采用全集成或混合集成形式。从电路设计角度看，核心关注点是器件之间的连接，因此具体集成形式并不会造成根本区别。本文用电子--光子异质融合集成电路（electronic-photonic heterogeneously-converging IC, EPHIC）表示电子与光子在电路层面的融合，而不限定具体集成形式。如图~1 所示，闭环反馈在传统 IC 和 EPHIC 中都具有重要作用~[2]。

为了支持高效的闭环 EPHIC 仿真，通用光子模型需要考虑以下特征：

\begin{enumerate}[label=\arabic*.]
  \item 模型必须能够描述光子器件随关键参数连续变化的特性。例如，在微环谐振器（MRR）的波长锁定中，需要获得 MRR 特性随温度的连续变化。
  \item 光子模型需要与已有电子模型兼容，从而支持大规模电子--光子系统设计。
  \item 闭环协同仿真同时包含光载波、电调制信号和热波动等不同时间尺度的物理量；这些量之间的频率失配会降低仿真效率。
  \item 模型应支持频域仿真，并能快速得到光子器件的频率响应。
  \item 模型还应考虑非线性效应、偏振和光信号双向传输。
\end{enumerate}

本文提出满足上述要求、用于闭环 EPHIC 仿真的通用 SPICE 解析光子模型，简称 EPHIC 模型~[3--5]，其结构如图~2 所示。第 II 节介绍现有光子模型，第 III 节介绍基础 EPHIC 模型，第 IV 节介绍复合 EPHIC 模型，第 V 节给出若干闭环电子--光子协同仿真实例，第 VI 节总结全文。

\noindent\textbf{图 1：}反馈在（a）传统 IC 和（b）EPHIC 中的作用~[2]。\par
\noindent\textbf{图 2：}用于闭环电子--光子协同仿真的 SPICE 光子模型框图。

\section{现有光子模型}

光子模型可采用散射参数模型、解析模型等不同实现原理；实现平台可以是电子设计自动化（EDA）、光子设计自动化（PDA）或其他平台；描述语言可以是 Verilog-A、SPICE 等。本节从这些角度讨论已有模型。

\subsection*{表 I：与现有光子模型的比较}

\begin{scriptsize}
\noindent
\begin{tabularx}{\linewidth}{>{\raggedright\arraybackslash}p{26mm}*{6}{>{\centering\arraybackslash}X}}
\toprule
项目 & [6] & [7] & [8] & [9] & [10] & EPHIC 模型\\
\midrule
平台 & Matlab & EDA & EDA & EDA & EDA & EDA\\
建模原理 & S 参数 & S 参数 & 解析 & 解析 & 解析 & 解析\\
模型类型 & 无源器件 & 无源器件 & 通用 & 通用 & 微环调制器 & 通用\\
描述语言 & Matlab & SPICE & Verilog-A & Verilog-A & SPICE & SPICE\\
偏振 & -- & -- & 支持 & -- & -- & 支持偏振管理\\
闭环仿真 & -- & -- & -- & 支持 & 支持 & 支持多维闭环\\
波长相关 & -- & 支持 & -- & -- & -- & 支持\\
非线性效应 & -- & -- & -- & -- & -- & 支持\\
基带等效 & 支持 & 支持 & -- & 支持 & -- & 支持\\
双向传输 & -- & 1 条总线 & 2 条总线 & 2 条总线 & -- & 1 条总线\\
模型复杂度（原文符号） & 2 & 2 & 3 & 3 & 1 & 5\\
\bottomrule
\end{tabularx}
\end{scriptsize}

散射参数模型通过复矢量拟合（CVF）算法拟合光子器件特性曲线~[6,11]。它对无源器件准确且有效；但有源光子器件在不同状态下具有不同的散射参数，因此必须用随状态变化的散射参数矩阵描述。闭环电子--光子协同仿真要求模型能够描述温度等参数连续变化时的物理特性，而散射参数模型难以直接支持这一要求。

紧凑解析光子模型~[9,12] 用物理方程描述器件。加入更多物理方程可以提高准确度，但也会降低仿真效率，因此解析模型需要在准确度和效率之间权衡。它的优势是能够用物理方程描述参数连续变化，因而适合闭环仿真。

EDA、PDA 和 Matlab 是常见的电子--光子协同仿真平台。Lumerical 使用散射参数和解析模型完成电路级仿真，而且器件级与电路级工具集成在同一平台中，便于从器件仿真中提取散射参数。Matlab 也支持解析模型~[13] 和散射参数模型~[6,11]，并擅长复杂矩阵运算。也可以用普通编程语言构建光子模型~[14]，但受功能和精度限制，这类方法不像 EDA/PDA 平台那样容易支持流片。EDA 已能进行大规模电子电路级仿真，因此在 EDA 中建立光子模型，是实现大规模电子--光子协同仿真的有效途径。

Verilog-A 和 SPICE 是常用的光子建模语言。已有工作用 Verilog-A 解析模型描述微环调制器（MRM）的光学和电学特性~[15,16]。但 Verilog-A 模型与 SPICE 电子模型协同仿真时，必须先编译为 AHDL，从而增加时间开销。文献~[12] 提出的基于频率啁啾的方法，比逐频点瞬态仿真更快地获得光子器件频率响应。文献~[10] 的 SPICE 兼容解析模型进一步提高了仿真效率，并能与电子电路良好结合；但该模型只面向 MRM，不够通用，不适合作为大规模光子建模方案。

文献~[7] 还实现了 SPICE 兼容的散射参数光子模型，能够准确描述器件特性，但仍只适用于无源器件，不能用于闭环电子--光子仿真。为了减小闭环电光电路中的频率失配，文献~[9] 的 Verilog-A 解析模型和文献~[7] 的 SPICE 兼容散射参数模型都采用基带等效方法，把光载频移到电调制频率附近。但当时尚未实现 SPICE 兼容的解析型基带等效模型。表 I 汇总了上述模型；原文认为，本文提出的 SPICE 兼容紧凑光子模型在所列项目上优于已有工作。

\transnote{表 I 是作者按自己选定的功能项做的定性比较。“复杂度”用对勾数量表达，并不是统一基准下测得的计算复杂度；因此应视为作者的方案定位，而不是独立性能排名。}

\section{基础 EPHIC 模型}

本节介绍波导、方向耦合器、连续波激光器、光纤--芯片耦合器、光电探测器、热效应、非线性效应和偏振等基础模型。

EDA 平台只能直接处理电子信号。SPICE 兼容光子模型的关键，是用数值等效的电子信号表示光信号。光信号包含相位、幅度、波长、偏振等多维信息；每一维可以用幅度和相位表示。但相位在 $2\pi$ 处存在跳变，因此本文不用“幅度--相位”，而用复光场的实部和虚部表示光信号~[17]。

一条电气总线可表示光信号的一个维度，其中两根信号线分别承载实部和虚部；多条总线可表示多个维度。与幅相表示~[18] 相比，实虚部表示无需增加总线数量就能支持多波长仿真，也便于用复乘（\texttt{cartmul}）和复加（\texttt{cartadd}）实现与频率无关的相移。由于电压和电流可以共存于同一总线，本文用电压表示光信号正向传播，用电流表示反向传播，如图~3。这样无需增加总线数量即可实现双向传输。

\noindent\textbf{图 3：}双向光传输的等效 SPICE 表示。CCCS 为电流控制电流源，VCVS 为电压控制电压源。

\subsection{波导}

波导是传输光信号的基础器件。传播的两个关键特征是延迟和损耗。波导末端 $z=L$ 的输出光场写为

\begin{equation}
E(L,t)=e^{\alpha L}E(0,t-t_d),
\tag{1}
\end{equation}

其中，$\alpha$ 是单位长度光损耗，$L$ 是波导长度，$t_d$ 是延迟。可用 VCVS 的 \texttt{gain} 属性描述损耗，用 \texttt{delay} 属性描述延迟。

\transnote{原文式 (1) 排成 $e^{\alpha L}$，但随后式 (4)、式 (12) 和式 (20) 都采用随正损耗系数衰减的 $e^{-\alpha L}$。实际实现前必须核对作者代码中的 $\alpha$ 符号定义；这里保留原文式 (1)，不擅自修正。}

为了提高效率，需要减小光信号与电信号之间的频率失配。基带等效模型~[6,9] 将光载频移到基带。光信号可写成~[9]

\begin{align}
E(z,t)&=e^{j\omega_R t}E_{\mathrm{shift}}(z,t)\notag\\
&=e^{-j\beta(\omega_R)z+j\omega_R t}
E_{\mathrm{shift}}\!\left(0,t-\frac{zn_g}{c}\right),
\tag{2}
\end{align}

其中，$\omega_R$ 是参考光信号角频率，$\beta$ 是波导传播常数，$n_g$ 是群折射率，$c$ 是光速，$E_{\mathrm{shift}}$ 是基带等效光信号。传播常数为

\begin{equation}
\beta=\frac{2\pi n_r}{\lambda},
\tag{3}
\end{equation}

其中 $n_r$ 是折射率，$\lambda$ 是波长。合并式 (1)--(3) 得

\begin{align}
E_{\mathrm{shift}}(L,t)
&=e^{-\alpha L-j\beta(\omega_R)L}
E_{\mathrm{shift}}\!\left(0,t-\frac{n_gL}{c}\right)\notag\\
&=e^{-\alpha L-j2\pi n_rL/\lambda}
E_{\mathrm{shift}}\!\left(0,t-\frac{n_gL}{c}\right).
\tag{4}
\end{align}

因此，基带等效光信号通过波导后，相对于原信号额外获得相移 $e^{-j2\pi n_rL/\lambda}$。后文模型中传输的信号均指基带等效光信号。

通用 SPICE 波导模型需要复数乘加。图~4 表明，\texttt{emcartmul} 和 \texttt{cartadd} 均可由 VCVS 实现。延迟后的输入信号 $E_{\mathrm{shift}}(0,t-t_d)$ 与基带转换因子 $e^{-j2\pi n_rL/\lambda}$ 共同输入 VCVS，生成 $E_{\mathrm{shift}}(L,t)$。对应波导模型见图~5。

\noindent\textbf{图 4：}（a）\texttt{cartmul} 与（b）\texttt{cartadd} 的等效 SPICE 表示。\par
\noindent\textbf{图 5：}波导的等效 SPICE 表示。

光子器件的频率响应可通过扫描波长 $\lambda$ 得到。与文献~[9,12] 相同，本文 SPICE 光子模型本质上是时域模型，支持瞬态仿真。由于电信号在这里仅作为光信号的数值等效量，它不能直接通过 SPICE 的电学 AC 分析得到光学频率响应。基于式 (4) 的基带模型，可把输入信号频率设为 0，并令 $\lambda$ 取不同值；EDA 工具用 DC 扫描即可自动计算每个波长的输出。

有效折射率随波长变化，原文采用

\begin{equation}
n_r(\lambda)=n_{r0}-(n_g-n_{r0})\frac{\lambda-\lambda_0}{\lambda_0},
\tag{5}
\end{equation}

其中 $\lambda_0$ 是参考波长，$n_g$ 是群折射率，$n_{r0}$ 是参考波长处的有效折射率。与基于啁啾信号的扫谱方法~[12] 和逐频点瞬态仿真~[9] 相比，EDA 中的 DC 扫描能更快得到光子器件频率响应。

为了支持闭环协同仿真，作者引入有源波导模型，以描述温度等因素引起的物性变化。从宏观上看，这些因素主要改变折射率、损耗和延迟。因此，可以把电光、热光和非线性效应归结为对折射率和损耗的修改，从而逐步扩展波导模型。

SPICE 没有直接支持连续可调传播延迟的通用器件，因此作者用 RC 网络近似延迟变化。当 $e^{j\omega t}$ 输入一阶 RC 网络时

\begin{equation}
\frac{V_{\mathrm{out}}}{V_{\mathrm{in}}}=\frac{1}{1+j\omega RC}.
\tag{6}
\end{equation}

当 $\omega RC\approx0$ 时

\begin{equation}
V_{\mathrm{out}}\approx
\frac{e^{j\omega t}}{\cos(\omega RC)+j\sin(\omega RC)}
=e^{j\omega(t-RC)},
\tag{7}
\end{equation}

即一阶 RC 网络可近似产生 $RC$ 的时延。可调延迟模块可用压控电阻或压控电容实现。由于光学频率特性通过 DC 扫描得到，RC 网络不会影响这种 DC 扫描结果。

不过，对长度超过约 10~mm 的波导，$\omega RC\approx0$ 不再成立。作者将延迟分为固定 DC 分量 $t_0$ 和变化 AC 分量 $t_1$：$t_0$ 用 VCVS 的固定延迟实现，$t_1$ 用可调 RC 网络实现。热光效应引起的变化量为

\begin{equation}
t_1=\frac{\Delta TL}{c}\frac{dn_g}{dT},
\tag{8}
\end{equation}

其中 $\Delta T$ 是波导温升，典型 $dn_g/dT=1.95\times10^{-4}\,\mathrm{K}^{-1}$~[19]。当温升为 10~K、波导短于 10~mm 时，一阶 RC 与标准延迟单元生成的信号基本一致：延迟相对误差约 0.56\%，幅度相对误差小于 0.1\%。更长波导需要更高阶 RC 网络。

\subsection{方向耦合器}

方向耦合器依据耦合模理论，把一个波导中的光场耦合到另一个波导~[20]：

\begin{align}
E_{\mathrm{out,top}}&=\gamma E_{\mathrm{in,top}}-j\kappa E_{\mathrm{in,bottom}},\notag\\
E_{\mathrm{out,bottom}}&=\gamma E_{\mathrm{in,bottom}}-j\kappa E_{\mathrm{in,top}}.
\tag{9}
\end{align}

$\gamma$ 是直通系数，$\kappa$ 是耦合系数。无损时

\begin{equation}
\gamma^2+\kappa^2=1.
\tag{10}
\end{equation}

式 (9) 可直接由 \texttt{cartmul} 和 \texttt{cartadd} 实现，如图~6。该方向耦合器本身忽略传播损耗；需要损耗时，可串接波导模型。

\noindent\textbf{图 6：}方向耦合器的等效 SPICE 表示。

\subsection{连续波激光器}

连续波（CW）激光器是光链路的激励源。本文把它简化成单频源，以便进行频率分析并降低模型复杂度。SPICE 兼容解析基带等效模型为

\begin{equation}
E_{\mathrm{out}}=E_0e^{j(\omega_0t+\phi)},
\tag{11}
\end{equation}

其中 $E_0$、$\phi$ 和 $\omega_0$ 分别是幅度、相位和移频后的角频率。可用正弦电压源 VSIN 的 \texttt{amplitude}、\texttt{frequency} 和 \texttt{phase} 属性描述。频率和幅度可调的激光源还可用 LC 谐振网络构造：压控电容调节谐振频率，VCVS 调节幅度。

\subsection{光纤--芯片耦合器}

光纤与片上波导的模场严重失配，因此高效率、宽带且兼容晶圆级测试的耦合仍有挑战。多数电路级仿真主要关注耦合效率与偏振敏感性。耦合效率可表示为额外损耗

\begin{equation}
E_{\mathrm{out}}=e^{-\alpha}E_{\mathrm{in}},
\tag{12}
\end{equation}

其中 $\alpha$ 表示耦合损耗项。

耦合器的偏振敏感性表现为不同偏振对应不同耦合效率。入射光的水平（$h$）和垂直（$v$）偏振分量，经过耦合器后分别转换为片上 TE 和 TM 模。对一维光栅耦合器，可把一个偏振分量的耦合效率设为 0，另一个设为测量或仿真值；对二维光栅耦合器和边缘耦合器，则分别给 $h$、$v$ 分量赋予测得或仿真的耦合效率，由此模拟偏振敏感性。

\subsection{光电探测器}

光电探测器（PD）把光强转换为电流。理想 PD 写为

\begin{equation}
I=\frac{1}{1+s/\omega_0}\,\mathcal{R}P+I_{\mathrm{dark}},
\tag{13}
\end{equation}

其中 $I$ 是光电流，$\mathcal{R}$ 是响应度，$P$ 是光功率，$\omega_0$ 是截止角频率，$I_{\mathrm{dark}}$ 是暗电流。光电流与光强近似线性，一阶低通滤波器描述 PD 的有限带宽。但响应度、带宽和寄生电容都会随偏置电压变化，因此实际模型需要依据代工厂工艺数据进一步拟合。式 (13) 中的响应度和暗电流可用受控源描述，带宽可用一阶 RC 描述，如图~7。

\noindent\textbf{图 7：}光电探测器的等效 SPICE 表示。

\subsection{热模型}

热光效应广泛用于硅光器件调谐，片上通常集成加热器。描述热调谐需要加热器模型、热传导模型和热光效应模型。已有工作用 RC 网络建立热模型~[21,22]；由于热光效应是闭环协同仿真的关键，原文仍给出详细过程。

电阻加热器产生的热功率为

\begin{equation}
P=I^2R,
\tag{14}
\end{equation}

其中 $I$ 是加热器电流，$R$ 是电阻。作者把温度视为“势”、热流视为“流”，从而把热阻映射为电阻、热容映射为电容，用 RC 网络描述一阶热传导，如图~8。该方法可以扩展到更复杂的热传递网络。最后，通过有效折射率的温度系数 $n_T$ 把温度变化耦合进光学模型。

\noindent\textbf{图 8：}加热器与一阶热传导模型的等效 SPICE 表示。

\subsection{非线性效应}

本文考虑光学 Kerr 效应~[23]、双光子吸收（TPA）~[24]、自由载流子吸收（FCA）和自由载流子色散（FCD）~[25]。

高光强通过 Kerr 效应改变折射率，近似为

\begin{equation}
\Delta n_{\mathrm{kerr}}=n_k I_{\mathrm{opt}},
\tag{15}
\end{equation}

其中 $n_k$ 是材料非线性折射率系数，$I_{\mathrm{opt}}$ 是光强。TPA 在两个光子的能量之和达到材料带隙时发生，其附加损耗近似为~[26]

\begin{equation}
\Delta\alpha_{\mathrm{TPA}}=\beta I_{\mathrm{opt}},
\tag{16}
\end{equation}

其中 $\beta$ 是 TPA 系数。TPA 产生自由载流子，继而引发 FCA 和 FCD。载流子密度满足~[26]

\begin{equation}
\frac{dN(t)}{dt}=\frac{\beta I^2(t)}{2h\nu}-\frac{N(t)}{\tau},
\tag{17}
\end{equation}

其中 $N(t)$ 是自由载流子密度，$h$ 是普朗克常数，$\nu$ 是光频率，$\tau$ 是复合寿命。载流子复合可用并联电阻表示，电阻值映射复合系数；光生载流子用压控电流源表示；载流子积累用并联电容表示，如图~9(b)。

在 1550~nm 附近，FCA、FCD 与载流子浓度的关系为~[25]

\begin{align}
\Delta\alpha_{\mathrm{FCA}}&=-8.5\times10^{-18}\Delta N_e
-6\times10^{-18}\Delta N_h,\notag\\
\Delta n_{\mathrm{FCD}}&=-8.8\times10^{-22}\Delta N_e
-8.5\times10^{-18}(\Delta N_h)^{0.8}.
\tag{18}
\end{align}

\transnote{式 (18) 的 FCA 项负号来自原文，但它与式 (4)、(16)、(19)、(20) 中把正的 $\alpha$ 或 $\Delta\alpha$ 作为衰减量的写法存在符号张力。实际实现前必须核对作者代码中 $\alpha$、$\Delta\alpha$ 和 $\Delta N$ 的符号约定，不能仅凭本式直接接入损耗模块。}

于是总附加损耗与折射率变化为

\begin{align}
\Delta\alpha&=\beta I_{\mathrm{opt}}+\Delta\alpha_{\mathrm{FCA}},\notag\\
\Delta n&=n_kI_{\mathrm{opt}}+\Delta n_{\mathrm{FCD}}.
\tag{19}
\end{align}

可用 VCCS 与串联电阻实现上述非线性模型，再把 $\Delta n$ 和 $\Delta\alpha$ 送入有源波导模型，见图~9(c)(d)。

自热是硅谐振结构中的重要非线性效应，来源于光吸收~[27]。FCA 和 TPA 引起的光功率损耗转化为热并加热波导；该效应更容易出现在高 $Q$ 的 MRR/MRM 中，并会改变器件工作状态。作者用波导输入、输出光功率之差估算自热：

\begin{align}
P_{\mathrm{sh}}&=\eta(I_{\mathrm{opt,in}}-I_{\mathrm{opt,out}})\notag\\
&=\eta\left[1-e^{-(\alpha_0+\Delta\alpha_{\mathrm{TPA}}+\Delta\alpha_{\mathrm{FCA}})L}\right]I_{\mathrm{opt,in}},
\tag{20}
\end{align}

其中 $\eta$ 是损失光功率转化为热的效率。可用以 $I_{\mathrm{opt,in}}$ 和 $I_{\mathrm{opt,out}}$ 为输入的 VCVS 及串联电阻描述，见图~9(a)。

\noindent\textbf{图 9：}（a）自热；（b）TPA 产生自由载流子；（c）FCA 与 TPA 引起的损耗变化；（d）FCD 与 Kerr 效应引起的有效折射率变化，对应的等效 SPICE 表示。

\subsection{偏振}

以往基于 Verilog-A 的光学模型只定义光的偏振状态~[8]。本文同时建模光的偏振状态和器件的偏振相关特性，使各种偏振态可以由不同偏振管理模型生成、操控、监测和提取。

模型用两条光学总线分别表示水平（$h$）和垂直（$v$）偏振；每条总线再用复光场的实部和虚部两根信号线表示。例如 $h$ 偏振由 $E_{h,\mathrm{Re}}(t)$ 和 $E_{h,\mathrm{Im}}(t)$ 表示。这种时域表示能够描述不同偏振分量实时的幅度与相位，并构造 Jones 向量

\begin{equation}
\begin{bmatrix}E_x\\E_y\end{bmatrix}
=\begin{bmatrix}\cos\alpha(t)\\\sin\alpha(t)e^{j\phi(t)}\end{bmatrix},
\tag{21}
\end{equation}

其中 $\alpha(t)$ 表示光场在 $h$、$v$ 分量间的幅度分配，$\phi(t)$ 是两偏振分量的相位差。它们可由实时光场得到：

\begin{align}
\alpha(t)&=\arctan\!\left(\frac{A_2(t)}{A_1(t)}\right),\notag\\
A_1(t)&=\sqrt{E_{h,\mathrm{Re}}^2(t)+E_{h,\mathrm{Im}}^2(t)},\notag\\
A_2(t)&=\sqrt{E_{v,\mathrm{Re}}^2(t)+E_{v,\mathrm{Im}}^2(t)},\notag\\
\phi(t)&=\arccos\!\left[
\frac{E_{h,\mathrm{Re}}E_{v,\mathrm{Re}}+E_{h,\mathrm{Im}}E_{v,\mathrm{Im}}}{A_1A_2}
\right].
\tag{22}
\end{align}

\transnote{式 (22) 忠实保留原文，但单独使用 $\arccos$ 只能返回 $[0,\pi]$，不能区分 $+\phi$ 与 $-\phi$，因此无法单独恢复完整相位象限和偏振旋向。完整实现通常应同时计算相应的点积与带符号叉积，并用 $\operatorname{atan2}$ 恢复相位差。}

Stokes 参数也可写成

\begin{align}
S_1&=\frac{A_1^2-A_2^2}{A_1^2+A_2^2},\notag\\
S_2&=\frac{2A_1A_2\cos\phi}{A_1^2+A_2^2},\notag\\
S_3&=\frac{2A_1A_2\sin\phi}{A_1^2+A_2^2}.
\tag{23}
\end{align}

由此可以提取和监测光的偏振态。原文声称，这是首次在电气平台上监测并操控光信号动态偏振态。

\section{复合 EPHIC 模型}

基础模型可组合为更复杂的光子器件。本节给出微环谐振器（MRR）、微环调制器（MRM）、马赫--曾德尔干涉仪（MZI）、马赫--曾德尔调制器（MZM）和偏振管理器件。

\subsection{微环谐振器}

最简单的 MRR 由直波导和环形波导组成。把方向耦合器与波导模型连接起来，即可构造 SPICE 兼容的 MRR。为了调谐谐振波长，模型还加入集成加热器和热模型，如图~10(a)。图~10(b) 中，本文 SPICE 模型的透射曲线与 Lumerical 和已有 Verilog-A 模型~[9] 一致。

仿真参数为：微环半径约 $10\,\mu\mathrm{m}$，功率耦合系数 $\kappa^2\approx0.25$，加热器电阻 $R_H\approx500\,\Omega$，有效折射率 $n_0\approx2$，群折射率 $n_g\approx3$，损耗 $\alpha\approx214.1978\,\mathrm{m}^{-1}$，参考波长 $\lambda_0\approx1550$~nm，环境温度 $T_0\approx300$~K。Lumerical、Verilog-A 和 SPICE 使用相同参数以公平比较；参数来自已有模型~[9] 和实验结果~[28]。除非另有说明，后续模型沿用这些参数。

加--下载型 MRR 在原结构旁增加第二条直波导，从环内耦出光场，适用于监测、开关和滤波。两个耦合区之间用两段直波导模型表示，并同样加入热模型，见图~11(a)；直通端和下载端透射曲线见图~11(b)。

级联 MRR 可提高带内平坦度、谱边沿陡峭度和消光比，是重要的光滤波方案。例如，二阶 MRR 可由三个耦合器和四段直波导构成，如图~12(a)，其直通端/下载端透射曲线见图~12(b)。

\noindent\textbf{图 10：}（a）MRR 原理图；（b）Lumerical、已有 Verilog-A 与本文 SPICE 模型得到的直通端透射曲线。\par
\noindent\textbf{图 11：}（a）加--下载型 MRR；（b）三种模型得到的直通端和下载端透射曲线。\par
\noindent\textbf{图 12：}（a）二阶 MRR；（b）三种模型得到的二阶 MRR 直通端/下载端透射曲线。

\subsection{微环调制器}

MRM 在环形波导中掺杂，通过载流子效应实现数据调制。本文建立载流子耗尽型 MRM 的 SPICE 兼容模型，如图~13(a)。有源波导的折射率、损耗和 PN 结寄生电容随调制电压变化，作者用四阶多项式拟合：

\begin{align}
n(V)&=n_0+n_1V+n_2V^2+n_3V^3+n_4V^4,\notag\\
\alpha(V)&=a_0+a_1V+a_2V^2+a_3V^3+a_4V^4,\notag\\
C_{pn}(V)&=c_0+c_1V+c_2V^2+c_3V^3+c_4V^4.
\tag{24}
\end{align}

其中 $V$ 是调制电压，$n$ 是有效折射率，$\alpha$ 是损耗，$C_{pn}$ 是 PN 结电容，各拟合系数来自文献~[28]。图~13(b)--(d) 给出 1~Gb/s NRZ 眼图、不同偏压下的透射曲线，以及不同光强下含自热效应的曲线。

加--下载型 MRM 可在加--下载型 MRR 基础上加入调制过程。MRM 还必须考虑波导非线性导致的自热；作者接入第 III 节的非线性与自热模型，仿真中出现了自热和双稳态，见图~13(d)。

\noindent\textbf{图 13：}（a）MRM 原理图；（b）1~Gb/s NRZ 光眼图；（c）不同偏置下的透射曲线；（d）不同光强下含自热的透射曲线，并比较 Lumerical、已有 Verilog-A 与本文 SPICE 模型。

\subsection{马赫--曾德尔干涉仪}

MZI 是硅光子学基础器件，通常由分光器、合光器、上臂和下臂组成。本文用两段不同长度的波导形成固定相位差，必要时再加入相移器。长度差 $\Delta L$ 引起

\begin{equation}
\Delta\phi=\frac{2\pi n\Delta L}{\lambda},
\tag{25}
\end{equation}

其中 $\lambda$ 是输入光波长，$n$ 是有效折射率。

其自由光谱范围为

\begin{equation}
\Delta\lambda_{\mathrm{FSR}}=\frac{\lambda^2}{n_g\Delta L}.
\tag{26}
\end{equation}

其中 $n_g$ 是群折射率。

模型用 Y 结构成分光器和合光器，用不同长度波导构成上下臂。片上 MZI 通常还用热相移器补偿工艺偏差或外部扰动；两臂相对相位随加热器电功率调整。与 Lumerical 的比较见图~14。

\noindent\textbf{图 14：}（a）MZI 原理图；（b）Lumerical 与 Cadence 中不同加热功率下的透射曲线。

\subsection{马赫--曾德尔调制器}

MZM 的上下臂通常包含热相移器、反向偏置 PN 相移器或正向偏置 PIN 相移器。正向偏置 PIN 结注入自由载流子，调制效率高、所需电压摆幅小；反向偏置 PN 结通过耗尽区载流子变化产生相移，寄生电容较小，可实现高速调制~[14]。不同相移器构成的 MZM 传输特性见图~15(b)(c)，其调制效率和损耗与其他平台的仿真一致。

\noindent\textbf{图 15：}（a）MZM 原理图；（b）正向偏置 PIN MZM；（c）反向偏置 PN MZM 的透射曲线。

\subsection{偏振管理器件}

光在光纤中传播时偏振态会变化。光纤引起的任意偏振变换可用酉矩阵表示：

\begin{equation}
T_{\mathrm{fiber}}=
\begin{bmatrix}e^{i\alpha}&0\\0&e^{-i\alpha}\end{bmatrix}
\begin{bmatrix}\cos\beta&-i\sin\beta\\-i\sin\beta&\cos\beta\end{bmatrix}
\begin{bmatrix}e^{i\gamma}&0\\0&e^{-i\gamma}\end{bmatrix},
\tag{27}
\end{equation}

其中 $\alpha$、$\beta$、$\gamma$ 是时变量。本文用三级 MZI 实现同一传递函数，如图~16(a)；各级相对相移对应这些时变量，上下臂分别对应 $h$、$v$ 偏振分量。相移器与 3~dB 耦合器可组合成偏振变化光纤模型。

偏振相关损耗可通过给 $h$、$v$ 分量设置不同耦合效率模拟；偏振模色散可通过给二者设置不同延迟模拟。片上主动偏振控制可克服偏振失配：如图~16(b)(c)，二维光栅耦合器，或带偏振分束旋转器的边缘耦合器，把光纤中的两个偏振转换成波导中的 TE 模，再由可调 MZI 合并，将任意输入偏振变换为固定 TE 模。外部电子反馈根据时变输入偏振相对调节 MZI 的两个相移器。

图~16(d)(e) 比较 Matlab 与 Cadence 结果。固定输入偏振为 TE/TM=1:1，扫描两个相移器的相移，即可得到各端口输出光功率变化。

\noindent\textbf{图 16：}（a）三级 MZI 构成的光纤模型；（b）基于二维光栅耦合器的主动偏振控制器；（c）基于边缘耦合器和偏振分束旋转器的控制器；（d）Matlab 与（e）Cadence 的传递函数。PS：相移器；TO：热光；WG：波导；PD：光电探测器；PSR：偏振分束旋转器。

\section{闭环电子--光子协同仿真}

作者用 MRR 波长锁定、MZM 偏置控制和偏振控制三个实例验证 EPHIC 模型。三组仿真均在 Cadence Virtuoso IC618 上进行，硬件为 Intel Xeon E5-2687W 与 128~GB 内存。由于 IHP 250~nm BiCMOS PDK 电路不能直接在 Lumerical 中设计，Lumerical 难以支持这里的完整闭环协同仿真，因此闭环部分主要比较 SPICE 兼容模型和已有 Verilog-A 模型。

\subsection{MRR 波长锁定}

作者基于 IHP 250~nm BiCMOS PDK 演示 MRR 闭环波长锁定。图~17(a) 中，跨阻放大器（TIA）采用简单并联反馈结构；控制器采用锁定最大值（lock-to-maximum, LTM）算法~[29]，目标是让探测光电流达到最大；功率级用功率 MOSFET 为集成加热器提供大电流。整个系统同时包含数字电路、模拟电路和光子器件。

图~17(b) 给出锁定过程。初始环境温度 $T_0=300$~K，微环温度 $T_{\mathrm{ring}}=300$~K，参考波长 $\lambda=1554$~nm。LTM 控制器逐渐增大功率级输出 $V_H$，使加热器升温，$T_{\mathrm{ring}}$ 和 PD 光电流 $I_{\mathrm{PD}}$ 增大。搜索完成后，$I_{\mathrm{PD}}$ 稳定在最大值，$T_{\mathrm{ring}}$ 稳定在对应值。给 $T_0$ 加入波动后，$V_H$ 能补偿扰动，使 $T_{\mathrm{ring}}$ 与 $I_{\mathrm{PD}}$ 基本不变。

Verilog-A 与本文 SPICE 模型结果一致。Lumerical 得到的 $T_{\mathrm{ring}}=300$~K 和 360~K 的透射曲线进一步表明，在 360~K、1554~nm 时微环透射达到目标极值；由 Lumerical 结果计算的光电流也与本文模型一致。作者还基于该模型实现并实验验证了含 MRR 的光子芯片，结果发表于文献~[30]。相同设置下，Verilog-A 闭环仿真耗时 38~s，EPHIC 模型耗时 27~s。

\noindent\textbf{图 17：}（a）MRR 波长锁定简化原理图；（b）锁定波形；（c）Lumerical 得到的 $T_{\mathrm{ring}}=300$~K、360~K 透射曲线。

\subsection{MZM 偏置控制}

MZM 偏置控制电路用于跟踪并补偿偏置点漂移~[31]，其集成化是大规模系统的重要环节。作者基于 IHP 250~nm BiCMOS PDK 演示闭环偏置控制。图~18(a) 的方案~[21] 将上下臂的一部分光耦合到 3~dB 耦合器，输出接平衡 PD 形成反馈。当 3~dB 耦合器输出达到最大或最小时，可建立正交相位条件。同时用 20~Gb/s PRBS-7 NRZ 信号调制 MZM。

TIA 把平衡 PD 电流转换为电压，经 ADC 数字化后送入微控制器。控制器估计偏置漂移，再通过 DAC 和热光相移器补偿。图~18(b) 中，$V_{\mathrm{crosstalk}}$ 加到另一热相移器上以模拟片上热扰动；反馈调节电压 $V_{\mathrm{heater}}$ 同时补偿初始相位误差并跟踪热波动，仿真结果与理论一致，反馈电流始终锁定在最小值。锁定与未锁定眼图对比显示偏置控制的必要性和有效性。相同设置下，Verilog-A 闭环仿真耗时 58~min，EPHIC 模型耗时 26~min。

\noindent\textbf{图 18：}（a）偏置控制方案简图；（b）仿真波形。

\subsection{偏振控制}

单模光纤中的偏振态会因扰动和制造误差随时间变化~[32]，片上主动偏振控制器可用于稳定偏振。作者基于 IHP 250~nm BiCMOS PDK 演示闭环主动偏振控制。图~19(a) 中，光学器件反馈端接 PD，电流经 TIA 和 ADC 送入控制器。控制器分析反馈波动，再通过 DAC 调节两个热相移器的相移，抑制波动，从而把任意输入偏振转换成固定片上偏振。

图~19(b) 中，通过改变幅度角 $\alpha$ 和相位角 $\phi$ 产生时变输入偏振；两个加热器电压跟踪输入偏振变化，并与理论结果一致。顶部 PD 电流和顶部 TIA 电压锁定在最大值，底部 PD 电流和底部 TIA 电压锁定在最小值。相同设置下，Verilog-A 闭环仿真耗时 2~h~12~min，EPHIC 模型耗时 1~h~56~min。

\noindent\textbf{图 19：}（a）主动偏振控制简图；（b）仿真波形。

\section{结论}

本文工作本质上是一个通用框架：用 SPICE 原生基元表示光信号，同时保留闭环电子--光子协同仿真所需的关键特征。EPHIC 模型支持较快的电子--光子协同仿真、双向光传输、光学非线性和多维信号等。其结果与已有 Verilog-A 和 Lumerical 模型一致，且框架还可继续扩展。

模型与现有 EDA 平台兼容，使 EPHIC 设计接近传统 IC 设计的便利程度。作者认为，这将推动光子学与电子学在电路层面进一步融合，并可能使 EPHIC 发展为与传统 IC 设计并列的独立研究方向。

\section*{参考文献}
\addcontentsline{toc}{section}{参考文献}
\small

\normalsize
\section*{附录：IEEE 原文页图}
\addcontentsline{toc}{section}{附录：IEEE 原文页图}
以下附入用户提供的 13 页原文，仅用于核对图、表、公式和译文位置。原文版权与使用限制仍按 IEEE 页面声明执行。

\includepdf[pages=-,pagecommand={\thispagestyle{empty}}]{ming_2024_ephic_models.pdf}

\end{document}
